\documentclass{article}
\usepackage{parskip}
\usepackage{amsmath}
\usepackage{amsthm}
\newtheorem{theorem}{Theorem}
\usepackage{amssymb}
\usepackage{pdflscape} 
\usepackage{fontspec}
% switch to a monospace font supporting more Unicode characters
\setmonofont{FreeMono}
\usepackage{minted}
\usepackage[most]{tcolorbox}
\tcbuselibrary{minted,skins}
% Define a breakable code box
\newcommand{\inputbreakableminted}[3]{%
  % #1 = language, #2 = file path, #3 = title
  \tcbinputlisting{%
    listing engine=minted,
    minted language=#1,
    listing file={#2},
    breakable,
    colback=gray!5!white,
    colframe=gray!75!black,
    title={#3},
    listing only,
    minted options={fontsize=\small, breaklines}
  }%
}

\definecolor{faintgray}{gray}{0.95} 
\newmintinline[lean]{lean4}{bgcolor=white}
\newminted[leancode]{lean4}{fontsize=\footnotesize}
\usemintedstyle{vs}  


\title{Lean Gauss Sum}
\author{Overleaf Support (Dan M)}
\date{September 2026}

\begin{document}
\maketitle


In this short example, we look at the Lean proof of the well known "Gauss sum" identity shown in Equation \ref{eq:gauss}.

\begin{theorem}[Sum of the First $n$ Natural Numbers]
For all natural numbers $n \in \mathbb{N}$, the sum of the first $n$ natural numbers satisfies the following identity:
\begin{equation}
    \sum_{i=0}^{i=n}i = \frac{n(n+1)}{2}
\end{equation}
\label{eq:gauss}
\end{theorem}

\begin{proof}
We proceed by mathematical induction on $n$. Let $S(n) = \sum_{i=1}^{n} i$. We wish to demonstrate that $2S(n) = n(n+1)$.

When $n = 0$, the sum is empty, yielding $S(0) = 0$. Evaluating both sides of the identity gives:
\[ 2(0) = 0(0 + 1) = 0 \]
Thus, the identity holds for $n=0$.

Assume the induction hypothesis holds for an arbitrary natural number $k$. That is, we assume:
\begin{equation}\label{eq:ih}
2S(k) = k(k + 1)
\end{equation}
We must now show that the identity holds for $k + 1$, meaning we want to prove that:
\[ 2S(k + 1) = (k + 1)(k + 2) \]

Using the recursive definition of summation, $S(n) = n + n-1 + n-2 + \ldots + 1$, we separate the final term $(k + 1)$ from the rest of the sum:
\[ S(k + 1) = (k + 1) + S(k) \]
Multiplying both sides by 2 yields:
\[ 2S(k + 1) = 2(k + 1) + 2S(k) \]
By applying our induction hypothesis from Equation~\eqref{eq:ih}, we substitute $2S(k)$ with $k(k + 1)$:
\[ 2S(k + 1) = 2(k + 1) + k(k + 1) \]
Factoring out the common term $(k + 1)$ from the right-hand side, we arrive at:
\[ 2S(k + 1) = (k + 1)(2 + k) = (k + 1)(k + 2) \]
This completes the inductive step. By the principle of mathematical induction, the identity holds for all natural numbers $n \in \mathbb{N}$.
\end{proof}

\section{The LEAN proof script}
The LEAN proof script for the Gauss summation identity is provided in Listing \ref{lst:gauss_code}. The proof script uses the LEAN language to express the approach taken by the human-oriented proof. 

\begin{listing}[!ht]
\inputminted[bgcolor=faintgray,breaklines]{lean4}{gauss.lean}
\caption{A Lean proof script of the Gauss summation identity}
\label{lst:gauss_code}
\end{listing}

A simpler version of the proof is possible by using LEAN \textbf{tactics}, which allows the proof to be written using imperative instructions that are expanded into full proofs by LEAN's elaborator. 

\begin{listing}[!ht]
\inputminted[bgcolor=faintgray,breaklines]{lean4}{gauss-tactic.lean}
\caption{A alternative LEAN proof script of the Gauss summation identity using tactics}
\label{lst:gauss_code2}
\end{listing}

\section{Compiling the proof script}

Running Listing \ref{lst:gauss_code} with LEAN will generate the formal proof that is provided at the end of this paper. A verification of the proof is provided by running the \texttt{\# print axioms} command against the generated proof. This inspects the compiled result and outputs an explicit list of  axioms used by the proof. Checking the output here shows that the proof relies only on standard axioms and does not use \texttt{sorry} or \texttt{unauthorized} axiom declarations that bypass real proof checking. The output from \texttt{\# print axioms} can be found in Listing \ref{lst:axioms}.

\begin{listing}[!ht]
\inputminted[bgcolor=faintgray,breaklines]{lean4}{leanchecker_output.txt}
\caption{Output from LEAN axiom checking}
\label{lst:axioms}
\end{listing}

The output from \ref{lst:axioms} shows that the proof relies exclusively on core  axioms: Propositional Extensionality (\texttt{propext}). Running Listing \ref{lst:gauss_code2} would add an additional dependency on Quotient Soundness (\texttt{Quot.sound}), and the non-constructive Axiom of Choice (\texttt{Classical.choice}).

%\begin{landscape}
%\section{Elaborated LEAN proof}
%\inputbreakableminted{lean4}{lean_output.txt}{LEAN Elaborated Proof}
%\end{landscape}

\begin{listing}[!ht]
\inputminted[bgcolor=faintgray,breaklines]{lean4}{lean_output.txt}
\caption{The elaborated Lean proof of the Gauss summation identity from Listing \ref{lst:gauss_code}}
\label{lst:elaboration_output}
\end{listing}
\end{document}
